\documentclass[12pt]{article}
\usepackage[a4paper,margin=1in]{geometry}
\usepackage{graphicx}
\usepackage[export]{adjustbox}
\usepackage{hyperref}
\usepackage{fontspec}
\usepackage[vietnamese]{babel}
\setmainfont{DejaVu Sans}
\usepackage{amsmath}
\usepackage{amssymb}
\usepackage{listings}

\begin{document}
BIÊN BẢN PHỎNG VẤN
VỊ TRÍ: QUẢN LÝ SẢN PHẨM CAO CẤP (GIẢI PHÁP AI)
ỨNG VIÊN
Alex Morgan
NGƯỜI PHỎNG VẤN
Sarah Jenkins (VP Prod)
NGÀY \& GIỜ
24 Tháng 10, 2026
VÒNG PHỎNG VẤN
Vòng 2 - Kỹ thuật \&
Lãnh đạo
Bối cảnh tổng quan: Phỏng vấn vòng hai đánh giá sự phù hợp chiến lược, nền tảng kỹ thuật trong phát triển sản phẩm AI và năng lực lãnh đạo của ứng viên cho vị trí dẫn dắt đội ngũ Nền tảng Đám mây Doanh nghiệp (Enterprise Cloud Platform).

PHẦN 1: THÔNG TIN NỀN TẢNG \& TẦM NHÌN
Sarah Jenkins
Chào mừng Alex. Cảm ơn bạn đã dành thời gian hôm nay. Để bắt đầu, bạn có thể tóm tắt ngắn gọn kinh nghiệm dẫn dắt các sáng kiến AI cho doanh nghiệp và điều gì đã thu hút bạn đến với vị trí này tại NexusTech không?
Alex Morgan
Cảm ơn Sarah. Rất hân hạnh được ở đây. Trong sáu năm qua tại CloudScale Inc., tôi đã quản lý các sản phẩm B2B SaaS, đặc biệt tập trung vào bộ phân tích học máy của chúng tôi. Tôi đã dẫn dắt quá trình chuyển đổi từ phân tích truyền thống sang các mô hình AI dự đoán, giúp tăng 35\% ARR. Tôi đã theo dõi sự mở rộng của NexusTech sang các giải pháp AI doanh nghiệp và rất ấn tượng với định hướng hạ tầng có thể mở rộng của quý công ty. Tôi thấy sự tương thích mạnh mẽ giữa kinh nghiệm tích hợp doanh nghiệp của tôi và lộ trình phát triển của công ty.
Sarah Jenkins
Đó là một thành tích rất ấn tượng. Bạn có thể đưa ra một ví dụ cụ thể về cách bạn ưu tiên các tính năng khi nợ kỹ thuật (technical debt) cạnh tranh trực tiếp với các yêu cầu cấp thiết từ phía khách hàng không?
Alex Morgan
Tất nhiên rồi. Năm ngoái, dịch vụ lập chỉ mục cốt lõi của chúng tôi gặp phải vấn đề nợ kỹ thuật từ hệ thống cũ, làm giảm tốc độ truy vấn gần 40\%. Cùng lúc đó, ba khách hàng doanh nghiệp lớn đang yêu cầu các tính năng báo cáo tùy chỉnh. Tôi đã áp dụng mô hình RICE có trọng số, dành cố định 20\% nguồn lực kỹ thuật mỗi sprint cho việc tái cấu trúc kiến trúc, đồng thời thương lượng giãn thời gian giao hàng cho các yêu cầu của khách hàng. Trong vòng hai quý, độ trễ truy vấn đã giảm một nửa và chúng tôi đã bàn giao 90\% bộ tính năng được yêu cầu mà không gây kiệt sức cho đội ngũ kỹ sư.

PHẦN 2: CHIẾN LƯỢC SẢN PHẨM \& GIẢI QUYẾT VẤN ĐỀ KỸ THUẬT
Sarah Jenkins
Bạn tiếp cận như thế nào khi thiết kế các tính năng dựa trên AI mà đầu ra của mô hình mang tính xác suất thay vì định tính (xác định)? Bạn quản lý kỳ vọng của người dùng và trải nghiệm UX xung quanh các lỗi như thế nào?
Alex Morgan
Thiết kế cho AI đòi hỏi sự chuyển dịch từ kết quả cố định sang các khoảng tin cậy. Đầu tiên, chúng tôi thiết kế sự minh bạch trực tiếp vào UI—hiển thị điểm tin cậy hoặc giải thích lý do tại sao một gợi ý được đưa ra. Thứ hai, chúng tôi đảm bảo các phương án dự phòng có sự can thiệp của con người (human-in-the-loop) diễn ra mượt mà. Nếu ngưỡng tin cậy giảm xuống dưới 80\%, UI sẽ gợi ý xem xét thủ công nhanh thay vì tự động ra quyết định một cách thầm lặng.

Đánh giá Tuyển dụng NexusTech • Ứng viên: Alex Morgan
Trang 1 / 2

PHẦN 2: CHIẾN LƯỢC \& GIẢI QUYẾT VẤN ĐỀ (TIẾP THEO)
Sarah Jenkins
Làm thế nào để bạn xác thực độ phù hợp sản phẩm - thị trường (product-market fit) trước khi đầu tư tài nguyên kỹ thuật lớn vào một mô hình AI mới?
Alex Morgan
Tôi kiên quyết ủng hộ phương pháp thử nghiệm 'Phù thủy xứ Oz' (Wizard of Oz) hoặc các bản mẫu thử nghiệm độ trung thực thấp (low-fidelity prototypes). Trước khi xây dựng các mô hình tùy chỉnh hoàn chỉnh, chúng tôi thử nghiệm nhu cầu bằng cách sử dụng API của bên thứ ba hoặc các quy trình bán thủ công để quan sát mức độ tương tác của người dùng. Chúng tôi đo lường thành công bằng các chỉ số định tính then chốt như tốc độ hoàn thành quy trình công việc và sự hài lòng định tính, thay vì các chỉ số hình thức ban đầu.

PHẦN 3: NĂNG LỰC LÃNH ĐẠO LIÊN PHÒNG BAN
Sarah Jenkins
Hãy mô tả một tình huống mà bộ phận kỹ thuật không đồng ý với chiến lược sản phẩm của bạn. Bạn đã làm thế nào để thống nhất các đội ngũ?
Alex Morgan
Trưởng nhóm kỹ thuật của chúng tôi muốn xây dựng lại toàn bộ đường ống dữ liệu (pipeline) bằng một framework mới, ước tính mất bốn tháng làm việc mà không đem lại

giá trị hướng tới khách hàng. Tôi đã tổ chức một buổi workshop chung, nơi chúng tôi đối chiếu các vấn đề nhức nhối của lập trình viên với dữ liệu rời bỏ của khách hàng. Chúng tôi đã thống nhất về một chiến lược chuyển đổi từng bước, giải quyết các nút thắt cổ chai nghiêm trọng nhất trước tiên trong khi vẫn tiếp tục phát hành các tính năng có tác động cao trong mỗi chu kỳ phát hành.

ĐÁNH GIÁ PHỎNG VẤN \& CHẤM ĐIỂM ỨNG VIÊN
LĨNH VỰC NĂNG LỰC
TÍN HIỆU / BẰNG CHỨNG QUAN SÁT ĐƯỢC
ĐÁNH GIÁ
Năng lực Kỹ thuật
Hiểu biết sâu sắc về quy trình làm việc AI, tích hợp API và các đánh đổi hệ thống.
4.5 /
5.0
Chiến lược Sản phẩm
Thể hiện khả năng sắp xếp thứ tự ưu tiên có cấu trúc bằng cách sử dụng các framework (RICE, MVP khám phá).
4.8 /
5.0
Quản lý Các bên liên quan
Thực tế trong việc tiếp cận các đánh đổi kỹ thuật và giao tiếp với khách hàng.
4.2 /
5.0
Lãnh đạo \& Văn hóa
Hợp tác, định hướng theo dữ liệu và thể hiện sự đồng cảm rõ ràng với các đội ngũ kỹ thuật.
4.6 /
5.0

Đề xuất của Người phỏng vấn: CHUYỂN TIẾP MẠNH MẼ
Alex đã thể hiện chuyên môn sâu sắc về quản lý sản phẩm B2B với sự nắm bắt xuất sắc các hạn chế của sản phẩm AI. Được đề xuất chuyển sang giai đoạn phỏng vấn cuối cùng với Phó Chủ tịch Kỹ thuật (VP of Engineering) và CPO.
Đánh giá Tuyển dụng NexusTech • Ứng viên: Alex Morgan
Trang 2 / 2
\end{document}
